% GENERATED FILE. Edit canonical lessons, then run npm run generate:books.
% canonical-source-hash: fnv1a64:9911e9c2ea96a328
% canonical-lessons: MW-C278-kahani, MW-C278-paheli, MW-C278-michno, MW-C278-putli, MW-C278-lukno, MW-R278-review-273-278

\chapter{Stories and Riddles}
\label{ch:mw-stories-and-riddles}
\hlchaptermodality{278}

\begin{chapteropening}
\textbf{By the end of this chapter:} \emph{I can say I like stories, riddles and hiding games.}

Complete the last lesson of chapter 278: I can say I like stories, riddles and hiding games.
\end{chapteropening}

\section[\texorpdfstring{\mw{कहाणी} (kahāṇī)}{कहाणी (kahāṇī)}]{\mw{कहाणी} (kahāṇī) — a story}

\label{lesson:MW-C278-kahani}



\begin{quote}
Before the new one: say the Marwadi for to sway, then the Marwadi for to be delighted.
\end{quote}

\subsection*{You'll want to know: \mw{कहाणी}}
\textbf{\mw{कहाणी}} — \emph{kahāṇī} — \textquotedblleft{}a story\textquotedblright{}.

\begin{quote}
\mw{दादी} \mw{रोज} \mw{कहाणी} \mw{सुणावै} \mw{है।}
\end{quote}

\begin{quote}
\emph{dādī roj kahāṇī suṇāvai hai.}
\end{quote}

\begin{quote}
Grandmother tells a story every day.
\end{quote}

A feminine noun.

\begin{cousinweb}
\textbf{\mw{कहाणी}} (\emph{kahāṇī}) is inherited from Sanskrit \emph{kathānaka}, \textquotedblleft{}a little tale\textquotedblright{}; CDIAL lists Marwari \emph{kā̃ṇi}, \textquotedblleft{}tale, proverb\textquotedblright{}, and Gujarati \emph{kahāṇī} (CDIAL 2705).
\end{cousinweb}

\begin{grammarlens}[title={What you've built}]
One more word for this chapter.
\end{grammarlens}

\subsection*{Your turn}
\begin{itemize}
\raggedright
  \item \emph{Say it:} \emph{kahāṇī}
  \item \emph{Say it:} \emph{kahāṇī}, once more
  \item \emph{Say it:} \emph{dādī roj kahāṇī suṇāvai hai.} — Grandmother tells a story every day.
  \item \emph{Recall:} say the Marwadi for to sway, then the Marwadi for to be delighted, then say \emph{kahāṇī} again
\end{itemize}

\begin{tcolorbox}[breakable,colback=teal!4,colframe=teal!35!black,title={Before you move on}]
What does \mw{कहाणी} mean? (\textquotedblleft{}A story\textquotedblright{}.) Say it once more.
\end{tcolorbox}
\section[\texorpdfstring{\mw{पहेली} (pahelī)}{पहेली (pahelī)}]{\mw{पहेली} (pahelī) — a riddle}

\label{lesson:MW-C278-paheli}



\begin{quote}
Before the new one: say the Marwadi for to be delighted, then the Marwadi for a story.
\end{quote}

\subsection*{You'll want to know: \mw{पहेली}}
\textbf{\mw{पहेली}} — \emph{pahelī} — \textquotedblleft{}a riddle\textquotedblright{}.

\begin{quote}
\mw{मास्टर} \mw{एक} \mw{पहेली} \mw{पूछै} \mw{है।}
\end{quote}

\begin{quote}
\emph{māsṭar ek pahelī pūchhai hai.}
\end{quote}

\begin{quote}
The teacher asks a riddle.
\end{quote}

A feminine noun.

\begin{cousinweb}
\textbf{\mw{पहेली}} (\emph{pahelī}) is inherited from Sanskrit \emph{prahelikā}, \textquotedblleft{}a riddle\textquotedblright{}; CDIAL gives no Marwari form in this entry, but lists Hindi \emph{pahelī}, \textquotedblleft{}riddle\textquotedblright{} (CDIAL 8911).
\end{cousinweb}

\begin{grammarlens}[title={What you've built}]
One more word for this chapter.
\end{grammarlens}

\subsection*{Your turn}
\begin{itemize}
\raggedright
  \item \emph{Say it:} \emph{pahelī}
  \item \emph{Say it:} \emph{pahelī}, once more
  \item \emph{Say it:} \emph{māsṭar ek pahelī pūchhai hai.} — The teacher asks a riddle.
  \item \emph{Recall:} say the Marwadi for to be delighted, then the Marwadi for a story, then say \emph{pahelī} again
\end{itemize}

\begin{tcolorbox}[breakable,colback=teal!4,colframe=teal!35!black,title={Before you move on}]
What does \mw{पहेली} mean? (\textquotedblleft{}A riddle\textquotedblright{}.) Say it once more.
\end{tcolorbox}
\section[\texorpdfstring{\mw{मीचणो} (mīchṇo)}{मीचणो (mīchṇo)}]{\mw{मीचणो} (mīchṇo) — to shut one's eyes}

\label{lesson:MW-C278-michno}



\begin{quote}
Before the new one: say the Marwadi for a story, then the Marwadi for a riddle.
\end{quote}

\subsection*{You'll want to know: \mw{मीचणो}}
\textbf{\mw{मीचणो}} — \emph{mīchṇo} — \textquotedblleft{}to shut one's eyes\textquotedblright{}.

\begin{quote}
\mw{छोरो} \mw{आंख} \mw{मीचै} \mw{है} \mw{अर} \mw{गिणै} \mw{है।}
\end{quote}

\begin{quote}
\emph{chhoro ā̃kh mīchai hai ar giṇai hai.}
\end{quote}

\begin{quote}
The boy shuts his eyes and counts.
\end{quote}

Past (I …): \mw{म्हैं} \mw{आंख} \mw{मीची}, \emph{mhaĩ ā̃kh mīchī}.

\begin{cousinweb}
\textbf{\mw{मीचणो}} (\emph{mīchṇo}) is inherited from a reconstructed \emph{*micc}, \textquotedblleft{}close the eyes, blink\textquotedblright{}; CDIAL gives no Marwari form in this entry, but lists Hindi \emph{mīcnā} and Gujarati \emph{micvũ}, both \textquotedblleft{}to close the eyes\textquotedblright{} (CDIAL 10118). In a hiding game the seeker shuts both eyes and counts.
\end{cousinweb}

\begin{grammarlens}[title={What you've built}]
One more word for this chapter.
\end{grammarlens}

\subsection*{Your turn}
\begin{itemize}
\raggedright
  \item \emph{Say it:} \emph{mīchṇo}
  \item \emph{Say it:} \emph{mīchṇo}, once more
  \item \emph{Say it:} \emph{chhoro ā̃kh mīchai hai ar giṇai hai.} — The boy shuts his eyes and counts.
  \item \emph{Recall:} say the Marwadi for a story, then the Marwadi for a riddle, then say \emph{mīchṇo} again
\end{itemize}

\begin{tcolorbox}[breakable,colback=teal!4,colframe=teal!35!black,title={Before you move on}]
What does \mw{मीचणो} mean? (\textquotedblleft{}To shut one's eyes\textquotedblright{}.) Say it once more.
\end{tcolorbox}
\section[\texorpdfstring{\mw{पुतळी} (putḷī)}{पुतळी (putḷī)}]{\mw{पुतळी} (putḷī) — a doll}

\label{lesson:MW-C278-putli}



\begin{quote}
Before the new one: say the Marwadi for a riddle, then the Marwadi for to shut one's eyes.
\end{quote}

\subsection*{You'll want to know: \mw{पुतळी}}
\textbf{\mw{पुतळी}} — \emph{putḷī} — \textquotedblleft{}a doll\textquotedblright{}.

\begin{quote}
\mw{म्हारी} \mw{पुतळी} \mw{नुवी} \mw{है।}
\end{quote}

\begin{quote}
\emph{mhārī putḷī nuvī hai.}
\end{quote}

\begin{quote}
My doll is new.
\end{quote}

A feminine noun.

\begin{cousinweb}
\textbf{\mw{पुतळी}} (\emph{putḷī}) is inherited from a reconstructed \emph{*putrala}, \textquotedblleft{}puppet\textquotedblright{}; CDIAL lists Old Marwari \emph{pūtaḷī}, \textquotedblleft{}image\textquotedblright{}, and Hindi \emph{putlī}, \textquotedblleft{}doll\textquotedblright{} (CDIAL 8269).
\end{cousinweb}

\begin{grammarlens}[title={What you've built}]
One more word for this chapter.
\end{grammarlens}

\subsection*{Your turn}
\begin{itemize}
\raggedright
  \item \emph{Say it:} \emph{putḷī}
  \item \emph{Say it:} \emph{putḷī}, once more
  \item \emph{Say it:} \emph{mhārī putḷī nuvī hai.} — My doll is new.
  \item \emph{Recall:} say the Marwadi for a riddle, then the Marwadi for to shut one's eyes, then say \emph{putḷī} again
\end{itemize}

\begin{tcolorbox}[breakable,colback=teal!4,colframe=teal!35!black,title={Before you move on}]
What does \mw{पुतळी} mean? (\textquotedblleft{}A doll\textquotedblright{}.) Say it once more.
\end{tcolorbox}
\section[\texorpdfstring{\mw{लुकणो} (lukṇo)}{लुकणो (lukṇo)}]{\mw{लुकणो} (lukṇo) — to hide (oneself)}

\label{lesson:MW-C278-lukno}



\begin{quote}
Before the new one: say the Marwadi for to shut one's eyes, then the Marwadi for a doll.
\end{quote}

\subsection*{You'll want to know: \mw{लुकणो}}
\textbf{\mw{लुकणो}} — \emph{lukṇo} — \textquotedblleft{}to hide (oneself)\textquotedblright{}.

\begin{quote}
\mw{छोरो} \mw{खाट} \mw{हेठै} \mw{लुकै} \mw{है।}
\end{quote}

\begin{quote}
\emph{chhoro khāṭ heṭhai lukai hai.}
\end{quote}

\begin{quote}
The boy hides under the bed.
\end{quote}

Past (he …): \mw{वो} \mw{लुक्यो}, \emph{vo lukyo}.

\begin{cousinweb}
\textbf{\mw{लुकणो}} (\emph{lukṇo}) is inherited from Sanskrit \emph{lupyate}, \textquotedblleft{}is lost, is removed\textquotedblright{}; CDIAL lists Old Marwari \emph{lukaï}, \textquotedblleft{}hides\textquotedblright{} (CDIAL 11083).
\end{cousinweb}

\begin{grammarlens}[title={What you've built}]
One more word for this chapter.
\end{grammarlens}

\subsection*{Your turn}
\begin{itemize}
\raggedright
  \item \emph{Say it:} \emph{lukṇo}
  \item \emph{Say it:} \emph{lukṇo}, once more
  \item \emph{Say it:} \emph{chhoro khāṭ heṭhai lukai hai.} — The boy hides under the bed.
  \item \emph{Recall:} say the Marwadi for to shut one's eyes, then the Marwadi for a doll, then say \emph{lukṇo} again
\end{itemize}

\begin{tcolorbox}[breakable,colback=teal!4,colframe=teal!35!black,title={Before you move on}]
What does \mw{लुकणो} mean? (\textquotedblleft{}To hide (oneself)\textquotedblright{}.) Say it once more.
\end{tcolorbox}
\section[(dialogue)]{Looking back: from Hush, a Secret to Stories and Riddles}

\label{lesson:MW-R278-review-273-278}



\begin{quote}
Say the Marwadi for silent and a gesture.
\end{quote}

\subsection*{Your turn}
Say these aloud:

\begin{itemize}
\raggedright
  \item silent, a gesture, a secret, to keep secret, a promise
  \item trouble, a cry for help, support, mercy, alms
  \item to praise, well done, a prize, credit for a deed, an opinion
  \item a flute, to play an instrument, a tune, a clap, to ring out
  \item a devotional song, a melody, rhythm, to sway, to be delighted
  \item a story, a riddle, to shut one's eyes, a doll, to hide
\end{itemize}

\begin{tcolorbox}[breakable,colback=teal!4,colframe=teal!35!black,title={Before you move on}]
Silent? (\textbf{\mw{चुप}}, \emph{chup}.) A gesture? (\textbf{\mw{इसारो}}, \emph{isāro}.) A secret? (\textbf{\mw{भेद}}, \emph{bhed}.) To keep secret? (\textbf{\mw{छिपावणो}}, \emph{chhipāvṇo}.) A promise? (\textbf{\mw{वचन}}, \emph{vachan}.) Trouble? (\textbf{\mw{तकलीफ}}, \emph{taklīph}.) A cry for help? (\textbf{\mw{पुकार}}, \emph{pukār}.) Support? (\textbf{\mw{सहारो}}, \emph{sahāro}.) Mercy? (\textbf{\mw{दया}}, \emph{dayā}.) Alms? (\textbf{\mw{भीख}}, \emph{bhīkh}.) To praise? (\textbf{\mw{सरावणो}}, \emph{sarāvṇo}.) Well done? (\textbf{\mw{साबास}}, \emph{sābās}.) A prize? (\textbf{\mw{इनाम}}, \emph{inām}.) Credit for a deed? (\textbf{\mw{जस}}, \emph{jas}.) An opinion? (\textbf{\mw{राय}}, \emph{rāy}.) A flute? (\textbf{\mw{बांसळी}}, \emph{bā̃sḷī}.) To play an instrument? (\textbf{\mw{बजावणो}}, \emph{bajāvṇo}.) A tune? (\textbf{\mw{सुर}}, \emph{sur}.) A clap? (\textbf{\mw{ताळी}}, \emph{tāḷī}.) To ring out? (\textbf{\mw{गूंजणो}}, \emph{gū̃jṇo}.) A devotional song? (\textbf{\mw{भजन}}, \emph{bhajan}.) A melody? (\textbf{\mw{राग}}, \emph{rāg}.) Rhythm? (\textbf{\mw{लय}}, \emph{lay}.) To sway? (\textbf{\mw{डोलणो}}, \emph{ḍolṇo}.) To be delighted? (\textbf{\mw{रीझणो}}, \emph{rījhṇo}.) A story? (\textbf{\mw{कहाणी}}, \emph{kahāṇī}.) A riddle? (\textbf{\mw{पहेली}}, \emph{pahelī}.) To shut one's eyes? (\textbf{\mw{मीचणो}}, \emph{mīchṇo}.) A doll? (\textbf{\mw{पुतळी}}, \emph{putḷī}.) To hide? (\textbf{\mw{लुकणो}}, \emph{lukṇo}.)
\end{tcolorbox}
